\section{Inheritance properties}

Este metodo surge como una extensión de los criterios de similitud de comporatmiento por permutación de etados y mediante reglas espejo. La idea core es que todas las reglas de todos los espacios posibles se encuentran conectadas, bajo el ejemplo de que la regla 0 en todos los sistemas siempre tiene el mismo comportamiento, por lo que no importa el numero de estados, el tamaño de la vecindad ni la dimensionalidad. Inherentemente surge la pregunta de ¿Existen otras reglas que tenga un comportamiento similar o heredado?

Esto surge considerando la regla 0 en el caso mas abstracto y sin sentido posible de un automata celular, un estado, un vecino (el mismo) y cero dimensiones (el unico objeto en su espacio), por lo que existe una unica regla que sera la regla 0, que se pude definir por una función constante $f(x) = 0$. Pero esta función constante eiste en todas las reglas 0 si aumentamos la dimensionalidad, la vecindad o los estados. Por lo que decimos que este es un comportamiento heredado. Lo propuesto son otras 2 a 3 tecnicas de similitud entre reglas pero a traves de definiciones (vecindad, dimensión, estados) las cuales son:

\begin{itemize}
	\item Reducción de estados
	\item reducción de vecindad
	\item Reducción de dimensionalidad
\end{itemize}

Por lo que decimos que 2 reglas tienen un comportamiento similar si al aplicar las tecnicas de redución y similitud (permutación de estados y espejo)  a ambas reglas, adquieren la misma regla. 

\subsection{Similitud entre reglas pertenecientes al mismo dominio}

Sean $\mathcal{A}_1 = (\mathcal{L}, \mathcal{S}, \mathcal{N}, \mathcal{R}_1)$ y $\mathcal{A}_2 = (\mathcal{L}, \mathcal{S}, \mathcal{N}, \mathcal{R}_2)$, dos reglas de automatas celulares diferentes existentes en el mismo dominio (dimensión, vecindad, estados), decimos que tienen un comportamiento similar si cumple alguna de las dos condiciones siguientes, donde $(x_1, \dots, x_m) \in \mathcal{S}^m$ denota los estados de la vecindad, $x_j = \sigma_{\vec{i} + \vec{v}_j}(t)$:

\begin{itemize}
	\item \textbf{Permutación de estados}: Si existe una biyección $\pi : \mathcal{S} \to \mathcal{S}$ tal que
	\[
	\mathcal{R}_2(x_1, \dots, x_m) = \pi\left( \mathcal{R}_1\left( \pi^{-1}(x_1), \dots, \pi^{-1}(x_m) \right) \right), \quad \forall (x_1, \dots, x_m) \in \mathcal{S}^m.
	\]
	
	Por ejemplo, para $\mathcal{S} = \{0, 1\}$, definimos $\pi(x) = 1 - x$, función que intercambia los dos estados.
	
	\item \textbf{Reglas espejo}:
	Si existe una reflexión $M$ que deja invariante a la vecindad ($M\mathcal{N} = \mathcal{N}$) tal que
	\[
	\mathcal{R}_2(x_1, \dots, x_m) = \mathcal{R}_1\left( x_{\rho(1)}, \dots, x_{\rho(m)} \right), \quad \forall (x_1, \dots, x_m) \in \mathcal{S}^m,
	\]
	donde $\rho$ es la permutación de índices definida por $\vec{v}_{\rho(j)} = M \vec{v}_j$. En una dimensión, $M\vec{v} = -\vec{v}$ y, ordenando la vecindad de izquierda a derecha, $\rho(j) = m + 1 - j$.
	
	Es decir, decimos que dos reglas son identicas si su comportamiento esta reflejado. 
\end{itemize}

TODO EJEMPLOS

\subsection{Reducción de dominio}

Las reducciones de dominio nos permiten disminuir la vecindad, la dimensionalidad o los estados de una regla, cada uno tiene su caso:

Sea $\mathcal{A} = (\mathcal{L}, \mathcal{S}, \mathcal{N}, \mathcal{R})$. La reducción de estados se considera cuando existen dos estados $a \neq b \in \mathcal{S}$ que cumplen estas dos condiciones:
\begin{enumerate}
	\item El estado $a$ no es producido. Es decir, $\mathcal{R}(x_1, \dots, x_m) \neq a, \; \forall (x_1, \dots, x_m) \in \mathcal{S}^m$.
	\item La transición de $a$ es idéntica a la de $b$. Es decir, sea $h : \mathcal{S} \to \mathcal{S} \setminus \{a\}$ con $h(a) = b$ y $h(x) = x$ para $x \neq a$:
	\[
	\mathcal{R}(x_1, \dots, x_m) = \mathcal{R}\left( h(x_1), \dots, h(x_m) \right), \quad \forall (x_1, \dots, x_m) \in \mathcal{S}^m.
	\]
\end{enumerate}
Entonces el autómata reducido es $\mathcal{A}' = (\mathcal{L}, \mathcal{S}', \mathcal{N}, \mathcal{R}')$ con $\mathcal{S}' = \mathcal{S} \setminus \{a\}$ y $\mathcal{R}' = \mathcal{R}|_{\mathcal{S}'^m}$.

El motivo de esta clasificación se resume a lo siguiente. Este estado $a$ nunca puede ser producido y ademas, su comportamiento coincide con el comportamiento del estado $b$, por lo que podemos y el comportamiento sera similar (con la condición de tener que sustituir o eliminar ese estado en la condición inicial de sistema ya que ese estado no pertenece nunca mas al nuevo automata celular definido por dicha reducción).

\subsection{Reducción de vecindad}

La reducción de vecindad ocurre cuando no importa el estado de un vecino, la salida sigue siendo la misma. Es decir, existe un índice $j$ tal que:
\[
\mathcal{R}(x_1, \dots, x_j, \dots, x_m) = \mathcal{R}(x_1, \dots, x_j', \dots, x_m), \quad \forall (x_1, \dots, x_m) \in \mathcal{S}^m, \; \forall x_j' \in \mathcal{S}.
\]
Entonces el autómata reducido es $\mathcal{A}' = (\mathcal{L}, \mathcal{S}, \mathcal{N}', \mathcal{R}')$ con $\mathcal{N}' = \mathcal{N} \setminus \{\vec{v}_j\}$ y $\mathcal{R}' : \mathcal{S}^{m-1} \to \mathcal{S}$ definida por
\[
\mathcal{R}'(x_1, \dots, x_{j-1}, x_{j+1}, \dots, x_m) = \mathcal{R}(x_1, \dots, x_m).
\]

Es decir, si en el vector de regla existe un vecino cuyo estado no afecta a la transición, entonces podemos aplicar una reducción de vecindad y eliminar este vecino. 

\subsection{Reducciones adicionales}

Adicionalmente se puede considerar la reducción de dimensionalidad, que al igual que los otros metodos, ocurriria cuando una dimensión no es utilizada (como establecer una regla que imite cualquier ECA en $\mathcal{L}^2$), esta no depende del eje adicional. Sin embargo, esta reducción no se considera para esta ivnestigación pues el objeto de estudio se limita a $\mathcal{L} ^1$.

CLAUDE: El siguiente parrafo tiene que definirse en el objeto de estudio de esta investigación (yo tengo la imagen)
Otra posible reducción podría ser la reducción de distancia vecinal, es decir, considerar unicamente a los vecinos por dirección y no por distancia. En palabras sencillas, los vectores de desplazamiento se definen ahora como una dirección y se consideran todos como vecinos inmediatos sin espacios entre celulas. No puede haber un vecino en -3 y otro en -1. Estas reglas de similitud no se estudian pero mi hipotesis es que es posible porque su comportamiento parece ser el mismo solo que bajo una distorción del espacio

TODO: incluir imagen de ejemplo de eso 

\subsection{Resultados de este metodo}

Con estas 4 nuevas clasificaciones podemos medir la similitud entre reglas de diferentes dominios. Aunque este método sigue siendo una clasificación binaria, es decir, solo nos dice si son o no son similares, ademas nos funciona como metodo de validación de validación para diferentes metodos, es decir, este dataset de relaciones nos ayuda a medir la exactitud de los metodos siguientes, pues ahora ya sabemos si son las mismas, por lo su indice de similitud bajo los metodos siguinetes deberia de ser cero.
